\documentclass[11pt]{article}
\usepackage[utf8]{inputenc}
\usepackage{amsmath,amssymb}
\usepackage{authblk}
\usepackage[margin=1in]{geometry}
\usepackage[hidelinks]{hyperref}
\usepackage{xurl}
\usepackage{setspace}

\title{RealQM, Standard Quantum Mechanics, and the Role of Relativity}
\author[1]{Claes Johnson}
\affil[1]{KTH Royal Institute of Technology, SE-100 44 Stockholm, Sweden \\
\texttt{claesjohnson@gmail.com}}
\date{\today \\[0.5em]
\small\textit{Written in respectful cooperation between Claes Johnson and Claude (Anthropic), with Claude
contributing a balanced, neutral view.}}

\begin{document}
\maketitle

\begin{abstract}
Standard quantum mechanics presents relativity as foundational to the physics of the atom: the electron requires
the Dirac equation, its spin and gyromagnetic factor $g=2$ are relativistic, fine structure is a relativistic
correction, light and matter join only through the relativistic field theory of QED, and inner electrons of heavy
atoms are said to \emph{move} at a large fraction of the speed of light. RealQM --- quantum mechanics as $N$ real
charge densities on non-overlapping domains in ordinary three-dimensional space, minimising a Coulomb energy with
no self-interaction --- is frankly non-relativistic, yet reproduces atomic energies across the periodic table,
covalent and hydrogen bonds, condensed phases, nuclear binding, and orbital magnetism. We ask what relativity was
actually doing, and answer that its role has been systematically overstated. The single phenomenon most cited as
forcing relativity onto the electron, $g=2$, is by L\'evy-Leblond's theorem a \emph{non-relativistic} geometric
fact --- the SU(2) double cover of the rotation group --- not a consequence of Lorentz invariance. The point
electron whose kinematics would demand relativity is replaced by a charge configuration obeying continuum
mechanics. And the heavy-atom ``relativistic speed'' story rests on a deeper confusion that RealQM exposes: in a
stationary state nothing moves ($\langle\mathbf p\rangle=0$, the density static), so standard QM's kinetic energy
$p^2/2m$ names a motion its own states forbid. RealQM reads the identical integral $\tfrac12\int|\nabla\psi|^2$
as what it physically is --- a localisation cost of static charge --- and the electron's ``velocity'' evaporates;
there is nothing to be relativistic. We are careful about the boundary: the Lorentz invariance of light, the
magnitude of fine structure, high-energy pair creation, and the QED anomaly are genuinely relativistic and are
coupled in classically, reconstructed, or left out of scope; and whether RealQM's static-charge functional needs
any high-$Z$ refinement is an open computational question, not a demonstrated need for relativity. Relativity, in
short, is \emph{sufficient} for much of what it is credited with and \emph{necessary} for far less.
\end{abstract}

\setstretch{1.06}

\section{Where standard quantum mechanics invokes relativity}

It is worth being specific about where the standard account reaches for relativity: the electron equation
(Dirac, 1928, taken as fundamental, with Schr\"odinger its low-velocity limit); spin and its factor $g=2$ (which
fall out of the Dirac equation and are called relativistic in origin); fine structure (the spin--orbit coupling
$\xi(r)\,\mathbf L\!\cdot\!\mathbf S$, a $v^2/c^2$ effect); antimatter and pair creation (with no non-relativistic
counterpart); and light (Lorentz invariant, coupled to matter through the relativistic field theory QED). The
first two of these, we will argue, are not what they seem; the last three mark the genuine boundary.

\section{RealQM builds the structure of matter without it}

The bulk of the physics of matter is the electrostatics of real charge and needs no relativity. Atomic total
energies across the periodic table, the covalent bond and its helium limit, hydride atomization energies, dimer
geometries, condensed phases, and even nuclear binding read as protons and electrons bound by Coulomb alone, all
follow from a non-relativistic variational principle on ordinary three-dimensional space. This is not an
approximation awaiting relativistic correction; for the structure, energetics, and chemistry of atoms it is the
physics itself, and relativity does no work.

\section{$g=2$ is geometric, not relativistic}

The single most-cited reason to believe the electron \emph{requires} relativity is $g=2$: since it drops out of
the Dirac equation, it is taken to be relativistic. It is not. L\'evy-Leblond (1967) derives $g=2$ from a
first-order, \emph{non-relativistic}, Galilean spinor equation: relativity is \emph{sufficient} to produce $g=2$
but not \emph{necessary}, because the factor of two does not come from Lorentz invariance in the first place. It
comes from the SU(2) double cover of the rotation group --- a spinor returns to itself only after $720^\circ$ ---
a geometric, scale-free fact about how the orientations of a real extended object compose, equally true in a
Galilean world. Writing the kinetic term with $\boldsymbol\sigma\!\cdot\!\mathbf p$ rather than $p^2$ produces
the $g=2$ Zeeman coupling automatically, with no relativity used. The phenomenon most responsible for the belief
that the electron is inescapably relativistic is, on inspection, non-relativistic geometry.

\section{Charges as configurations, not relativistic objects}

Underlying the standard picture is the electron as a point object with intrinsic rest mass, a thing whose very
kinematics is relativistic. RealQM reads elementary charges instead as \emph{configurations} of charge density in
space, held by the same Coulomb energetics as everything else, their apparent mass and rigidity emergent
properties of the configuration rather than primitive relativistic attributes. There is then no point electron
whose motion demands Lorentz kinematics; there is an extended distribution obeying continuum mechanics, and the
``particle'' is a stable pattern in it. The need for relativity as the kinematics of point particles dissolves
with the point particle.

\section{Heavy atoms, and a kinetic energy without motion}

The most-cited evidence that atoms are relativistic comes from the heavy elements, and within standard quantum
chemistry it is real: non-relativistic calculations predict silvery gold and solid mercury, relativistic ones
predict the observed yellow gold and liquid mercury, the gap growing as $(Z\alpha)^2$. That is a fact about those
calculations. How it should be read is another matter, and less settled than the standard story assumes.

The standard telling ascribes the effect to inner electrons \emph{moving} at a large fraction of $c$. This is
incoherent for a stationary bound state, where nothing moves: $\langle\mathbf p\rangle=0$ and the density is
static. Standard QM's kinetic energy
\[
\Big\langle \tfrac{p^2}{2m}\Big\rangle = \tfrac12\int|\nabla\psi|^2\,d^3x
\]
is called a \emph{kinetic} energy while nothing is kinetic --- a formalism without a physical referent for the
motion it names. The heavy-atom ``relativistic velocity'' reifies the formal quantity $\langle p^2\rangle$, which
is a \emph{spread}, into a \emph{speed} the theory's own dynamics forbid.

RealQM gives the identical integral its honest meaning. In RealQM $\tfrac12\int|\nabla\psi|^2$ is not the energy
of a moving particle but a \emph{localisation cost} of a static charge density --- a strain-like elastic energy
penalising rapid spatial variation, the same energy that stabilises matter against collapse. With that reading
the electron never moves, at any $Z$, and the ``relativistic speed'' of heavy-atom electrons simply evaporates:
there is no velocity to be relativistic.

Two consequences follow, and both matter. First, whether RealQM's localisation functional needs a higher-order
term at extreme compression (the heavy-atom regime) is an \emph{open and untested} question --- not a settled
deficiency, and certainly not a call for relativistic motion. The heavy-atom differential properties (the gold
$5d$--$6s$ gap, the mercury cohesion) have not been computed in full spatial detail in RealQM. Second, the
few-percent deviations in RealQM's heavy-atom \emph{total} energies are the accumulation of a simplified,
spherically-symmetric computation of a large multi-shell number, not a relativistic shortfall. It is a bad habit
of modern physics to charge every residual imprecision to ``relativistic effects''; RealQM declines it. If a
full RealQM computation should turn out to need a high-$Z$ correction, it will be a statement about the
\emph{energy} of a still charge, never about a motion the theory's own stationary states forbid.

\section{The honest boundary}

What genuinely requires relativity, and is therefore coupled in classically, reconstructed, or left out of
scope? Light is Lorentz invariant, but a real charge density couples directly to a classical Maxwell field as a
source, with no configuration-space amplitude to promote to a quantum field --- the elaborate machinery of QED,
needed in the standard theory precisely because its matter lives on $3N$-space, is not required to join real
charge to real light. High-energy pair creation and the parton regime are genuinely relativistic and outside a
Coulomb continuum theory of bound charge. The measured $g=2.00232\ldots$ is a QED radiative effect; RealQM
reaches the geometric $g=2$, not its correction. Spacetime itself RealQM does not presuppose: the quiet atom is
timeless, and relativity is reconstructed constructively from observation rather than assumed a fixed backdrop.

\section{Conclusion}

The role of relativity in the physics of matter has been systematically overstated. For the structure,
energetics, chemistry, and orbital magnetism of atoms, molecules, and nuclei, relativity does no work, and RealQM
computes these without it. $g=2$ is geometry, not relativity; the point electron is replaced by a charge
configuration; and the heavy-atom ``relativistic speed'' is exposed as an artefact of reading $p^2/2m$ as a
motion that stationary states do not contain. What is genuinely relativistic --- the invariance of light, the
magnitude of fine structure, high-energy processes, the QED anomaly --- is coupled in classically, left honestly
out of scope, or reconstructed. Relativity is \emph{sufficient} for much of what it is credited with and
\emph{necessary} for far less; and the atom's electrons, in RealQM, are always still.

\begin{thebibliography}{9}
\bibitem{RealQM} C.~Johnson, \emph{Quantum Mechanics as Multiphase 3D Continuum Mechanics}, manuscript, 2026;
\url{https://claes542.github.io/RealMolecule/gallery.html}.
\bibitem{Spinor} C.~Johnson, \emph{The Spinor Residue: $g=2$ and Stern--Gerlach without Relativity}, manuscript,
2026.
\bibitem{Charges} C.~Johnson, \emph{Elementary Charges as Configurations, not Objects --- Why RealQM Needs No
Relativity}, manuscript, 2026.
\bibitem{Time} C.~Johnson, \emph{Real Time Dependence in RealQM}, manuscript, 2026.
\bibitem{ManyMinds} C.~Johnson, \emph{Many-Minds Relativity}, Body and Soul, Vol.~VII, 2026.
\bibitem{LevyLeblond} J.-M.~L\'evy-Leblond, \emph{Nonrelativistic particles and wave equations}, Commun.\ Math.\
Phys.\ \textbf{6}, 286 (1967).
\end{thebibliography}

\end{document}
